\begin{textsample}{POS Dim 1 – pi – Score 62.00 – pi/pi\_em\_43.txt}  \label{ex:f1_pos_052}
\textbf{Competências} e Habilidades com a \textbf{Competência} Geral 2 da BNCC, no que se refere ao exercício da curiosidade intelectual que utiliza o conhecimento para investigar, refletir e criar soluções em diferentes situações. ( CE02 ) Articular conhecimentos matemáticos ao propor e/ou participar de ações para investigar desafios do mundo \textbf{contemporâneo} e tomar decisões éticas e socialmente responsáveis, com base na \textbf{análise} de \textbf{problemas} de urgência social, como os voltados a situações de saúde, sustentabilidade, das implicações da \textbf{tecnologia} no mundo do trabalho, entre outros, recorrendo a conceitos, procedimentos e linguagens próprios da Matemática.
Esta \textbf{competência} está voltada para a investigação de ações possíveis de serem realizadas a partir de \textbf{análises} com base em elementos da matemática, tendo em \textbf{vista} aspectos éticos e socialmente responsáveis. \textbf{Note} que a ênfase dessa \textbf{competência} está na intervenção na realidade.
Essas ações de intervenção poderão fornecer um excelente contexto para aprender novos conceitos matemáticos ao mesmo tempo em que permitem aos \textbf{alunos} vivenciar a \textbf{aplicabilidade} da matemática.
A \textbf{competência} 2 coloca o estudante como personagem atuante em sua comunidade local e no mundo \textbf{globalizado}.
As ações de propor e participar fazem referência à \textbf{capacidade} de ser parte de algo, compartilhar saberes com o outro e colaborar conjuntamente para a produção de algo.
Destaca -se também o papel da investigação por parte do estudante, o que pressupõe a observação dos desafios presentes em sua comunidade local/global, a elaboração de hipóteses que as descrevam, o tratamento dos dados associados à situação envolvida, a \textbf{análise} dos resultados obtidos e, por fim, a tomada de decisão a partir das conclusões obtidas.
Ao desenvolver essa \textbf{competência}, pode -se afirmar que o estudante avança em relação ao entendimento de que os Projetos de Vida não são apenas no âmbito profissional, mas também nas dimensões pessoal e social/cidadã.
Está associada à \textbf{Competência} Geral 7 da BNCC que permite ao estudante um \textbf{posicionamento} ético em relação ao cuidado de si mesmo, dos outros e do planeta, bem como com a \textbf{Competência} Geral 10 que trata do agir pessoal e coletivo com tomada de decisões, a partir de princípios éticos, democráticos, inclusivos, sustentáveis e solidários. ( CE03 ) Utilizar estratégias, conceitos e procedimentos matemáticos, em seus campos – Aritmética, Álgebra, Grandezas e Medidas, Geometria, Probabilidade e Estatística – para interpretar, construir modelos e resolver \textbf{problemas} em diversos contextos, analisando a plausibilidade dos resultados e a adequação das soluções propostas, de modo a construir argumentação consistente.
A C3 prevê a utilização de estratégias, conceitos, definições e procedimentos matemáticos para construir modelos e resolver \textbf{problemas}.
Uma vez construído o modelo e resolvido o \textbf{problema}, é preciso se debruçar sobre essa solução de modo a analisar a adequação e veracidade do elaborado para convencer a si, aos colegas e ao professor de que sua conclusão é correta. \textbf{Convém} também observar que no processo de resolução de \textbf{problemas} o estudante deve ser motivado a : questionar … formular … testar e validar hipóteses … buscar contraexemplos … modelar situações … \textbf{verificar} adequação da resposta.
E, como consequência, desenvolver linguagens e construir formas de pensar que o levem a refletir e agir de \textbf{maneira} \textbf{crítica}.
É importante \textbf{frisar} que a referida \textbf{competência} não se \textbf{restringe} apenas à resolução de \textbf{problemas}, mas também trata de sua elaboração.
Isso revela uma concepção da resolução de \textbf{problemas} além da \textbf{mera} aplicação de um conjunto de regras.
Outro grande destaque refere -se à modelagem matemática como a construção de modelos matemáticos que sirvam para generalizar ideias ou para descrever situações semelhantes.
Essa \textbf{competência} tem estreita relação com a \textbf{Competência} Geral 2 da BNCC, no sentido da \textbf{capacidade} de formular e resolver \textbf{problemas}, e com a \textbf{Competência} Geral 4, que reforça a importância de saber utilizar as diferentes linguagens para expressar ideias e informações para a comunicação mútua. ( CE04 ) Compreender e utilizar, com flexibilidade e fluidez, diferentes registros de representação matemáticos ( algébrico, geométrico, estatístico, computacional, entre outros ), na busca de solução e comunicação de resultados de \textbf{problemas}, de modo a favorecer a construção e o desenvolvimento do \textbf{raciocínio} matemático.
Nesta \textbf{competência} há a proposição do uso flexível dos diferentes registros de representação da linguagem matemática.
O domínio de tais registros amplia a \textbf{capacidade} de \textbf{análise} das situações representadas e \textbf{aprimora} \textbf{tanto} os aspectos de formulação como a comunicação de resultados.
As aprendizagens previstas nessa \textbf{competência} contribuem ainda para melhorar a \textbf{capacidade} de \textbf{argumentar} e identificar \textbf{raciocínios} falsos.
A \textbf{competência} 4 complementa as demais no sentido de que utilizar, interpretar e resolver situações-problema se faz pela comunicação das ideias dos estudantes por meio da linguagem matemática.
Transitar entre os diversos tipos de representações ( simbólica, algébrica, gráfica, textual etc. ) permite a \textbf{compreensão} mais profunda dos conceitos e ideias da matemática.
A representação de uma mesma situação de diferentes formas estabelece conexões que possibilitam resolver \textbf{problemas} matemáticos usando estratégias diversas.
Além disso, a \textbf{capacidade} de elaborar modelos matemáticos para expressar situações \textbf{implica} e revela a aprendizagem, além de potencializar o \textbf{letramento} matemático.
Essa \textbf{competência} está relacionada ao desenvolvimento das \textbf{Competências} Gerais 4 e 5 da BNCC, uma vez que a linguagem utilizada de modo flexível permite expressar ideias e informações que facilitam o entendimento e ampliar o repertório de formas de expressão, inclusive a \textbf{digital} com espaço para \textbf{autoria} pessoal e criatividade do estudante. ( CE05 ) Investigar e estabelecer conjecturas a respeito de diferentes conceitos e propriedades matemáticas, empregando recursos e estratégias como observação de padrões, experimentações e \textbf{tecnologias} \textbf{digitais}, identificando a necessidade, ou não, de uma demonstração cada vez mais formal na validação das referidas conjecturas.
A \textbf{competência} 5 prevê que os \textbf{alunos} potencializem seus processos de abstração e generalização na elaboração de conjecturas.
Busquem recursos que permitam uma validação mais formal da conjectura elaborada ou sua refutação.
Os \textbf{alunos} devem utilizar o \textbf{raciocínio} indutivo \textbf{baseado} na observação de padrões e em regularidades.
No entanto, eles deverão reconhecer que o \textbf{raciocínio} indutivo e as verificações meramente empíricas não são apropriadas para validar propriedades, podendo levar, algumas vezes, a conclusões equivocadas.
Desse modo, é necessário reconhecerem que as conjecturas formuladas devem ser justificadas por meio de resultados matemáticos já demonstrados.
Assim, a \textbf{competência} 5 tem como objetivo principal que os estudantes se apropriem da forma de pensar matemática, como \textbf{ciência} com uma forma específica de validar suas conclusões pelo \textbf{raciocínio} lógico-dedutivo.
Não se trata de trazer para o Ensino Médio a Matemática formal dedutiva, mas de permitir que os jovens percebam a diferença entre uma dedução originária da observação empírica e uma dedução formal.
É importante também \textbf{verificar} que essa \textbf{competência} e suas habilidades não se desenvolvem em separado das demais ; ela é um \textbf{foco} a mais de atenção para o ensino em \textbf{termos} de formação dos estudantes, de modo que identifiquem a Matemática diferenciada das demais Ciências.
As habilidades para essa \textbf{competência} \textbf{demandam} que o estudante vivencie a investigação, a formulação de hipóteses e a tentativa de validação de suas hipóteses.
De \textbf{certa} forma, a proposta é que o estudante do Ensino Médio possa conhecer parte do processo de construção da Matemática, tal qual aconteceu ao longo da história, fruto do pensamento de muitos em diferentes culturas.
Um ponto de atenção está no fato de que algumas das habilidades escolhidas pela BNCC ( 2018 ) para essa \textbf{competência} \textbf{remetem} a conteúdos muito específicos, de pouca \textbf{aplicabilidade} e de difícil \textbf{contextualização}, mas que, no entanto, favorecem a investigação e a formulação de hipóteses antes de que os estudantes conheçam os conceitos ou a \textbf{teoria} subjacente a esses conteúdos específicos.
As habilidades propostas para essa \textbf{competência} possuem níveis diferentes de complexidade cognitiva, desde a identificação de uma propriedade até a investigação completa com dedução de uma regra ou procedimento.
Essa \textbf{competência} se relaciona com as \textbf{Competências} Gerais 2, 4, 5 e 7 da BNCC, uma vez que há o incentivo ao exercício da curiosidade intelectual na investigação, neste caso, com maior centralidade no conhecimento matemático.
A linguagem e os recursos \textbf{digitais} são ferramentas básicas e essenciais para facilitar a observação de regularidades, expressar ideias e construir \textbf{argumentos} com base em fatos.
O conjunto das \textbf{competências} específicas da área de Matemática estão articuladas entre si, possibilitando aos estudantes identificarem a importância do conhecimento matemático em seu desenvolvimento pessoal, social e profissional.
SÍNTESE DAS \textbf{COMPETÊNCIAS} Aplicar conhecimentos Construir \textbf{Argumentos} Comunicar em Matemática Formalizar/Demonstrar PROCESSOS PARA APRENDIZAGEM Investigação Construção de Modelos Resolução de \textbf{Problemas} As habilidades selecionadas para a \textbf{matriz} curricular de matemática As habilidades previstas na \textbf{matriz} curricular de Matemática para a etapa do Ensino Médio na BNCC se articulam com as \textbf{competências} gerais e específicas da área e também com as aprendizagens já previstas para o ensino fundamental.
No entanto, como \textbf{explicitado} no documento da Base : > [ … ] Por sua vez, embora cada habilidade esteja associada a determinada \textbf{competência}, isso não significa que ela não contribua para o desenvolvimento de outras. [ … ] as habilidades são apresentadas sem \textbf{indicação} de \textbf{seriação}.
Essa decisão permite flexibilizar a definição anual dos currículos e propostas pedagógicas de cada escola. ( p. 530 ) Essa característica permite escolher as habilidades livres da pressão da \textbf{seriação} e, por um critério de decidir por sua centralidade na formação integral, pela importância dela na constituição da área, pelas conexões que favorecem na área ou entre áreas e, ainda, pela possibilidade de seu desenvolvimento em períodos reduzidos de escolaridade, tais como programas de aceleração da aprendizagem ou cursos de Educação de Jovens e Adultos ( EJA ).
Na \textbf{sequência}, apresentaremos o \textbf{quadro} destinado à Matriz Curricular de Matemática indicada por série, de modo a possibilitar melhor adequação ao trabalho pedagógico dos professores, ao que inclui -se a carga horária prevista.
Além disso, são especificados as \textbf{Competências} Específicas ( CE ) da área de Matemática ( CE01, CE02, CE03, CE04, CE05 ) e as \textbf{Competências} Gerais da BNCC as quais se relacionam ; os objetos do conhecimento e objetivos de aprendizagem para cada habilidade, segundo a \textbf{competência} específica ( CE ) estabelecida.
Para melhor \textbf{compreensão}, ressaltamos que as habilidades são o pressuposto para a \textbf{indicação} dos objetos do conhecimento e objetivos de aprendizagem, o que permitiu elaborar a \textbf{matriz} por série.
Por isso, é possível que habilidades diferentes incorram em objetos do conhecimento semelhantes ou iguais.
No entanto, o que será determinante na \textbf{abordagem} deste objeto do conhecimento pelo professor é, unicamente, a complexidade \textbf{exigida} pela habilidade.
Assim, por exemplo, o \textbf{tópico} de “ juros \textbf{simples} ” está inserido na 1ª e 3ª séries, porém com habilidades e objetivos de aprendizagem distintos, na primeira são observados aspectos relativos à definição e cálculo e na outra, aplicações deste conceito em objetivos mais \textbf{complexos} da matemática financeira ; o mesmo acontece com outros objetos do conhecimento especificados na \textbf{matriz}.
O currículo proposto nesse formato possibilita que o \textbf{aluno} possa ir do conhecimento geral ao conhecimento especializado naturalmente.
Isso se deve a um modelo de aprendizagem contínua que evita que os conceitos caiam facilmente no esquecimento.
Como mencionado na \textbf{introdução} deste documento, o desenvolvimento integral, as dez \textbf{competências} gerais e os princípios \textbf{metodológicos} integradores são a principal forma de integrar o desenvolvimento de ações educacionais de diversos tipos a partir da \textbf{matriz} curricular.
No entanto, é possível observar situações específicas nas quais a Matemática se relaciona com as demais áreas, possibilitando a construção de conhecimentos de modo contextualizado pelos estudantes.
A Matemática pode ser entendida como uma linguagem, sendo parte dos processos de letrar os \textbf{alunos} em diferentes sentidos, formando leitores para múltiplas formas de comunicação, \textbf{aproximando} -se assim da área de Linguagens.
As Artes Gráficas e a Música são outras promissoras formas de integração da Matemática com a área de Linguagens.
Conhecimentos de geometria ampliam a percepção de espaço e das formas de sua representação, fornecendo referenciais para as representações planas nos desenhos, mapas e ferramentas de localização espacial.
Tais representações são úteis a \textbf{Ciências} da Natureza e Ciências Humanas.
As habilidades relacionadas à probabilidade e estatística, em especial aquelas ligadas a coletar, organizar, descrever e analisar dados, realizar inferências e fazer predições com base em amostras de população, são aplicações da Matemática em questões do mundo real que tiveram um crescimento muito grande e se tornaram instrumentos importantes também para as Linguagens e as Ciências Humanas.
A articulação entre Matemática e Ciências da Natureza se dá pela leitura, pela escrita e pela resolução de \textbf{problemas}, nos quais a Matemática que sintetizam as investigações, a generalização de resultados encontrados e a possibilidade de utilizá -los na resolução de \textbf{problemas} semelhantes.
Além disso, projetos temáticos se constituirão cenários \textbf{propícios} à integração da Matemática com as demais áreas de conhecimento.
Outro aspecto a ser considerado na implementação da \textbf{matriz} curricular é o planejamento.
Planejar a utilização das metodologias integradoras de forma mais adequada para o ensino de cada objeto de conhecimento ou habilidade não pode ser feito no momento da \textbf{aula}, uma vez que é preciso considerar os \textbf{alunos}, os recursos disponíveis, o tempo para o ensino e o acompanhamento das aprendizagens \textbf{esperadas}, sendo que a avaliação é \textbf{central} nesse processo.
A avaliação em seu caráter \textbf{diagnóstico} é feita sempre que for necessário ouvir e observar os estudantes, para permitir que expressem ou registrem o que sabem, ou pensam que sabem, sobre um conceito, um \textbf{problema}, uma observação sistemática etc.
A avaliação deve cumprir também seu papel de acompanhamento do processo de aprendizagem, quando o professor observa, registra, solicita que o estudante fale ou represente como pensou, como fez, ou o que não entendeu.
Na perspectiva de uma avaliação formativa, é preciso considerar as limitações dos instrumentos \textbf{prova}, listas de exercícios ou pesquisas sem alinhamento com um projeto, muito utilizados apenas para gerar notas.
Isso significa \textbf{qualificar} mais o instrumento prova, adotando diferentes modalidades para o emprego desse recurso de avaliação, e inserir outras formas de avaliar validadas pelo registro e pela clareza do que se está avaliando.
Um elemento da \textbf{matriz} que é muito \textbf{relevante} para o planejamento e a avaliação está traduzido em “ exemplos de objetivos de aprendizagem ”.
Eles indicam aquilo que se \textbf{espera} de aprendizagem para cada conjunto de \textbf{competências} e habilidades e, podem ser \textbf{parâmetros} para acompanhar ensino e aprendizagem, fazer ajustes na caminhada e planejar formas de conseguir que todos aprendam.

% matched lemmas: abordagem, aluno, análise, aplicabilidade, aprimorar, aproximar, argumentar, argumento, aula, autoria, basear, capacidade, central, certo, ciência, competência, complexo, compreensão, contemporâneo, contextualização, convir, crítico, demandar, diagnóstico, digital, esperar, exigir, explicitar, foco, frisar, globalizado, implicar, indicação, introdução, letramento, maneira, matriz, mero, metodológico, notar, parâmetro, posicionamento, problema, propício, prova, quadro, qualificar, raciocínio, relevante, remeter, restringir, sequência, seriação, simples, tanto, tecnologia, teoria, termos, tópico, verificar, vista
\end{textsample}
